%% Dokumentenklasse (Koma Script) -----------------------------------------
\documentclass[%
   %draft,     % Entwurfsstadium
   final,      % fertiges Dokument
	 % --- Paper Settings ---
   paper=a4,% [Todo: add alternatives]
   paper=portrait, % landscape
   pagesize=auto, % driver
   % --- Base Font Size ---
   fontsize=12pt,%
	 % --- Koma Script Version ---
   version=last, %
 ]{scrbook} % Classes: scrartcl, scrreprt, scrbook

% Encoding der Dateien (sonst funktionieren Umlaute nicht)
\usepackage{listings}
\usepackage[utf8]{inputenc}

%%% Preambel
%%% === Textbody ==============================================================
\KOMAoptions{%
   DIV=15,% (Size of Text Body, higher values = greater textbody)
   % DIV=calc % (also areaset/classic/current/default/last) 
   % -> after setting of spacing necessary!   
   BCOR=5mm% (Bindekorrektur)
}%
%%% === Headings ==============================================================
\KOMAoptions{%
   %%%% headings
	% headings=small,  % Small Font Size, thin spacing above and below
   % headings=normal, % Medium Font Size, medium spacing above and below
   headings=big, % Big Font Size, large spacing above and below
   %
   headings=noappendixprefix, % chapter in appendix as in body text
	headings=nochapterprefix,  % no prefix at chapters
   % headings=appendixprefix,   % inverse of 'noappendixprefix'
   % headings=chapterprefix,    % inverse of 'nochapterprefix'
   % headings=openany,   % Chapters start at any side
   % headings=openleft,  % Chapters start at left side
   headings=openright, % Chapters start at right side      
   %%% Add/Dont/Auto Dot behind section numbers 
   %%% (see DUDEN as reference)
   % numbers=autoenddot
   % numbers=enddot
   numbers=noenddot
   % secnumdepth=3 % depth of sections numbering (???)
}%
\setcounter{secnumdepth}{3}
%%% === Page Layout ===========================================================
\KOMAoptions{% (most options are for package typearea)
   twoside=true, % two side layout (alternating margins, standard in books)
   % twoside=false, % single side layout 
   % twoside=semi,  % two side layout (non alternating margins!)
   %
   twocolumn=false, % (true)
   %
   headinclude=false,%
   footinclude=false,%
   mpinclude=true,%      
   %
   headlines=2.1,%
	% headheight=2em,%
   headsepline=true,%
   footsepline=false,%
   cleardoublepage=empty %plain, headings
}%
%%% === Paragraph Separation ==================================================
\KOMAoptions{%
	% parskip=relative, % change indentation according to fontsize (recommanded)
   parskip=absolute, % do not change indentation according to fontsize
   parskip=false    % indentation of 1em
   % parskip=true   % parksip of 1 line - with free space in last line of 1em
   % parskip=full-  % parksip of 1 line - no adjustment
   % parskip=full+  % parksip of 1 line - with free space in last line of 1/4
   % parskip=full*  % parksip of 1 line - with free space in last line of 1/3
   % parskip=half   % parksip of 1/2 line - with free space in last line of 1em
   % parskip=half-  % parksip of 1/2 line - no adjustment
   % parskip=half+  % parksip of 1/2 line - with free space in last line of 1/3
   % parskip=half*  % parksip of 1/2 line - with free space in last line of 1em
}%
%%% === Table of Contents =====================================================
\setcounter{tocdepth}{3} % Depth of TOC Display
\KOMAoptions{%
   %%% Setting of 'Style' and 'Content' of TOC
   % toc=left, %
   toc=indented,%
   %
   toc=bib,
   %toc=nobib,
   % toc=bibnumbered,
   %
	% toc=index,%
   toc=noindex,
   %
   % toc=listof,
   toc=nolistof
   % toc=listofnumbered,
   %   
}%  
%%% === Lists of figures, tables etc. =========================================
\KOMAoptions{%
   %%% Setting of 'Style' and 'Content' of Lists 
   %%% (figures, tables etc)
	% --- General List Style ---
   listof=left, % tabular styles
   listof=indented, % hierarchical style
   % --- chapter highlighting ---
   % listof=chapterentry, % ??? Chapter starts are marked in figure/table
   % listof=chaptergapline, % New chapter starts are marked by a gap 
      		  			   	 % of a single line
	listof=chaptergapsmall, % New chapter starts are marked by a gap 
   	    					   % of a smallsingle line
   % listof=nochaptergap, % No Gap between chapters
   %
   % listof=leveldown, % lists are moved one level down ???
   % --- Appearance of Lists in TOC
    listof=notoc % Lists are not part of the TOC
   %listof=totoc % add Lists to TOC without number
   % listof=totocnumbered, % add Lists to TOC with number
}%  
%%% === Bibliography ==========================================================
%% Setting of 'Style' and 'Content' of Bibliography
\KOMAoptions{%
	% bibliography=oldstyle,%
   bibliography=openstyle,%
   bibliography=nottotoc % Bibliography is not part of the TOC
   % bibliography=totocnumbered, % add Bibliography to TOC with number
   %bibliography=totoc % add Bibliography to TOC without number
}%
%%% === Index =================================================================
%% Setting of 'Style' and 'Content' of Index in TOC
\KOMAoptions{%
   index=nottotoc % index is not part of the TOC
	% index=totoc, % add index to TOC without number
}%
%%% === Titlepage =============================================================
\KOMAoptions{%
   titlepage=true %
   %titlepage=false %
}%
%%% === Miscellaneous =========================================================
\KOMAoptions{% 	
   footnotes=multiple% nomultiple
   %open=any,%
   %open=left,%
   %open=right,%
   %chapterprefix=false,%
   %appendixprefix=false,%
   %chapteratlists=10pt,% entry
}%

\gdef\thetitle{}
\gdef\thethesis{}
\gdef\theauthor{}
\gdef\theabstract{}
\gdef\theabgabedatum{}
\gdef\thegeburtsort{}
\gdef\thegeburtsdatum{}
\gdef\thematrikel{}
\gdef\thesupervisor{}
\gdef\thezusammenfassung{}
%
\def\title#1{\gdef\@title{#1}\gdef\thetitle{#1}}
\def\@title{\@latex@error{No \noexpand\title given}\@ehc}
\def\thesis#1{\gdef\thethesis{#1}}
\def\author#1{\gdef\@author{#1}\gdef\theauthor{#1}}
\def\@author{\@latex@warning@no@line{No \noexpand\author given}}
\def\abgabedatum#1{\gdef\theabgabedatum{#1}}
\def\geburtsort#1{\gdef\thegeburtsort{#1}}
\def\geburtsdatum#1{\gdef\thegeburtsdatum{#1}}
\def\abstract#1{\gdef\theabstract{#1}}
\def\matrikel#1{\gdef\thematrikel{#1}}
\def\supervisor#1{\gdef\thesupervisor{#1}}
\def\zusammenfassung#1{\gdef\thezusammenfassung{#1}}
\input{preambel/preambel}
%
%%%% Neue Befehle
\def\bzgl{bzgl.\ }
\def\dh{d.\,h.\ }
\def\eV{e.\,V.\ }
\def\engl{engl.:\ }
\def\zB{z.\,B.\ }
\def\oae{o.\,\"{a}.\ }
\def\uU{u.\,U.\ }
\def\ua{u.\,a.\ }
\def\ia{i.\,a.\ }
\def\og{o.\,g.\ }
\def\bzw{bzw.\ }
\def\vgl{vgl.\ }
\def\etc{etc.\ }
\def\evtl{evtl.\ }
\def\ca{ca.\ }
\def\sog{sog.\ }
\def\insbes{insbes.\ }
\def\usw{usw.\ }

%--------------------------- References ------------------------------
\def\Eq#1{\eqref{#1}}
\def\Sec#1{Abschnitt~\ref{#1}}
\def\SubSec#1{Subsection~\ref{#1}}
\def\Ex#1{Beispiel~\ref{#1}}
\def\Chap#1{Kapitel~\ref{#1}}
\def\Part#1{Teil~\ref{#1}}
\def\Fig#1{Abbildung~\ref{#1}}
\def\Tab#1{Tabelle~\ref{#1}}
\def\Theo#1{Satz~\ref{#1}}
\def\Def#1{Definition~\ref{#1}}
\def\Alg#1{Algorithmus~\ref{#1}}
\def\Cor#1{Korollar~\ref{#1}}
\def\Lem#1{Lemma~\ref{#1}}
\def\Ax#1{Axiom~\ref{#1}}
\def\Page#1{Seite~\pageref{#1}}

%-------------------------- Debugging ----------------------------------
\def\work#1{#1}
\def\draft#1{}
\def\final#1{}
\def\todo#1{\work{\comment{TODO: #1}}}

%%%%%%%%%%%%%%%%%%%%%%%%%%%%%%%%%%%%%%%%%%%%%%%%%%%%%%%%%%%%%%
%%
%%   Identification
%%
%%%%%%%%%%%%%%%%%%%%%%%%%%%%%%%%%%%%%%%%%%%%%%%%%%%%%%%%%%%%%%
%\NeedsTeXFormat{LaTeX2e}

%%%%%%%%%%%%%%%%%%%%%%%%%%%%%%%%%%%%%%%%%%%%%%%%%%%%%%%%%%%%%%
%%
%%   Preliminary declarations
%%
%%%%%%%%%%%%%%%%%%%%%%%%%%%%%%%%%%%%%%%%%%%%%%%%%%%%%%%%%%%%%%
%\RequirePackage{bm} % correct bold symbols, like \bm
%\RequirePackage{amsmath}
%\RequirePackage{amsfonts}
%\RequirePackage{amssymb}
%\RequirePackage{mathrsfs}
%\RequirePackage{dsfont} % double stroke symbols like IR
%\RequirePackage[single]{accents} % own accents 
%\RequirePackage{exscale} % Scaling of Math Symbols?

%%%%%%%%%%%%%%%%%%%%%%%%%%%%%%%%%%%%%%%%%%%%%%%%%%%%%%%%%%%%%%
%%
%%   Options
%%
%%%%%%%%%%%%%%%%%%%%%%%%%%%%%%%%%%%%%%%%%%%%%%%%%%%%%%%%%%%%%%



%%%%%%%%%%%%%%%%%%%%%%%%%%%%%%%%%%%%%%%%%%%%%%%%%%%%%%%%%%%%%%
%%
%%   More declarations
%%
%%%%%%%%%%%%%%%%%%%%%%%%%%%%%%%%%%%%%%%%%%%%%%%%%%%%%%%%%%%%%%


%------------------------------------------------------------
% Mathematical Formulae
%------------------------------------------------------------
\newcommand{\rv}[1]{\ensuremath{{\boldsymbol{#1}}}}
\renewcommand{\vec}[1]{\ensuremath{{\underline{#1}}}}
\newcommand{\rvec}[1]{\ensuremath{{\boldsymbol{\underline{#1}}}}}
\newcommand{\mat}[1]{{\ensuremath{{\mathbf{#1}}}}}

\newcommand{\Menge}[1]{{\mathcal{#1}}}
\newcommand{\SquareRoot}[1]{{\mbox{\bfseries\sffamily #1}}}
\newcommand{\MatrixMenge}[2]{{\mathcal{M}}_{#1#2}}
\newcommand{\intndim}{\ds\int_{-\infty}^\infty \dots \int_{-\infty}^\infty}
\newcommand{\Eye}{\mat{I}}

% operators (in alphabetical order):
\renewcommand{\arg}{\operatorname*{arg}}
%\DeclareMathOperator*{\argmin}{arg\,min}
%\DeclareMathOperator*{\argmax}{arg\,max}
\newcommand{\argmin}[1]{\mathop{\mathrm{arg}\; \min_{#1}}}
\newcommand{\argmax}[1]{\mathop{\mathrm{arg}\; \max_{#1}}}
\DeclareMathOperator{\col}{col}
\newcommand{\dd}{\operatorname{d}\!}
%\def\d{\ensuremath{\operatorname{d}}} % for \int f(x)\d x or \d y / \d y
%\renewcommand{\d}{\ensuremath{\text{d}}}
\DeclareMathOperator{\diag}{diag}
\DeclareMathOperator{\Cov}{Cov}
\DeclareMathOperator{\erf}{erf}
\DeclareMathOperator{\E}{E}
\DeclareMathOperator{\Proj}{proj}
\renewcommand{\Pr}{\operatorname{Pr}} 
\DeclareMathOperator{\rang}{rang}
\DeclareMathOperator{\rank}{rank}
\DeclareMathOperator{\jj}{j}% complex unit
\DeclareMathOperator{\rect}{rect}
\def\T{^\mathrm{T}} % transpose
\DeclareMathOperator{\trace}{trace}
\DeclareMathOperator*{\Var}{Var} 
\renewcommand{\arctan}{\operatorname{arctan}}
\newcommand{\atann}{\operatorname{arctan2}}
%\DeclareMathOperator{\atan}{atan}
\def\atan2{\operatorname{atan2}}
\DeclareMathOperator{\pot}{pot}

\newcommand{\MatrixTyp}[2]{\Menge{M}_{#1, #2}}
\newcommand{\FromTo}[2]{$#1 = 1,\, \ldots,\, #2$}
\newcommand{\LMInd}{{\operatorname{LM}}}

%---------------------------------------------------------------------
% Abbreviations
%---------------------------------------------------------------------

\def\({\left(}
\def\){\right)}
\def\[{\left[}
\def\]{\right]}

\newcommand{\ds}{\displaystyle}
\newcommand{\ns}{\textstyle}

\newcommand{\NewB}{\ensuremath{\mathds{R}}}
\newcommand{\NewC}{\ensuremath{\mathds{C}}}
\newcommand{\NewN}{\ensuremath{\mathds{N}}}
\newcommand{\NewQ}{\ensuremath{\mathds{Q}}}
\newcommand{\NewR}{\ensuremath{\mathds{R}}}
\newcommand{\NewZ}{\ensuremath{\mathds{Z}}}
\newcommand{\IR}{\NewR}
\newcommand{\IN}{\NewN}
\newcommand{\IZ}{\NewZ}
\newcommand{\IC}{\NewC}

\newcommand{\Gauss}{{\mathcal{N}}}

\setcounter{MaxMatrixCols}{20} %Fuer ams umgebung


%--------------------------------------------------------------------------
%Einheiten mit korrekten Leerzeichen
%Usage: \Meter{\micro}{10} results $10\mu m$ 
\newcommand{\Meter}[2][]{\unit{#2}#1\meter} 
\newcommand{\Volt}[2][]{\unit{#2}#1\volt} 
\newcommand{\Hertz}[2][]{\unit{#2}#1\hertz} 
\newcommand{\Second}[2][]{\unit{#2}#1\second} 
\newcommand{\Pascal}[2][]{\unit{#2}#1\pascal} 
\newcommand{\Euro}[2][]{\unit{#2}#1\euro}

\newcommand{\nm}[1] {\Meter[\nano]{#1}}
\newcommand{\mum}[1]{\Meter[\micro]{#1}}
\newcommand{\mm}[1] {\Meter[\milli]{#1}}
\newcommand{\cm}[1] {\Meter[\centi]{#1}}
\newcommand{\mV}[1] {\Volt[\milli]{#1}}
\newcommand{\ms}[1] {\Second[\milli]{#1}}
\newcommand{\KHz}[1] {\Hertz[\kilo]{#1}}

%--------------------------- Operators ---------------------------
\def\TJ{(TJ)}
\def\TkJ{(T^kJ)}
\def\TJp{(T_\pi J)}
\def\TkJp{(T_\pi^k J)}
\def\vTJ{T  \vec J}
\def\vTkJ{T^k\vec J}
\def\vTJp{T_\pi \vec J}
\def\vTkJp{T_\pi^k \vec J}

%------------------------ Random variables ---------------------------
\def\v{\rv v}
\def\w{\rv w}
\def\x{\rv x}
\def\y{\rv y}
\def\z{\rv z}

%---------------------------- Vectors --------------------------------
\def\vdel{\vec \delta}
\def\veta{\vec \eta}
\def\vOmega{\vec \Omega}
\def\vsig{\vec \sigma}
\def\va{\vec a}
\def\vb{\vec b}
\def\vc{\vec c}
\def\ve{\vec e}
\def\vg{\vec g}
\def\vh{\vec h}
\def\vmu{\vec \mu}
\def\vn{\vec n}
\def\vpi{\vec \pi}
\def\vu{\vec u}
\def\vv{\vec v}
\def\vw{\vec w}
\def\vx{\vec x}
\def\vy{\vec y}
\def\vz{\vec z}
\def\vxi{\vec \xi}
\def\vB{\vec B}
\def\vF{\vec F}
\def\vJ{\vec J}
\def\vM{\vec M}
\def\vP{\vec P}
\def\vX{\vec X}

%------------------------- Random vectors ----------------------------
\def\rvn{\rvec n}
\def\rvv{\rvec v}
\def\rvw{\rvec w}
\def\rvx{\rvec x}
\def\rvy{\rvec y}
\def\rvz{\rvec z}
\def\rvhz{\rvec{\hat z}}
\def\rvxi{\rvec \xi}
\def\rvX{\rvec X}

%---------------------------- Matrices -------------------------------
\def\A{\mat A}
\def\B{\mat B}
\def\C{\mat C}
\def\D{\mat D}
\def\mE{\mat E}
\def\F{\mat F}
\def\H{\mat H}
\def\I{\mat I}
\def\K{\mat K}
\def\L{\mat L}
\def\M{\mat M}
\def\P{\mat P}
\def\Q{\mat Q}
\def\R{\mat R}
\def\V{\mat V}
\def\W{\mat W}

%------------------------------ Mean ---------------------------------
\def\ha{\hat a}
\def\hv{\hat v}
\def\hw{\hat w}
\def\hx{\hat x}
\def\hy{\hat y}
\def\hz{\hat z}
\def\hva{\hat{\vec a}}
\def\hvv{\hat{\vec v}}
\def\hvw{\hat{\vec w}}
\def\hvx{\hat{\vec x}}
\def\hvy{\hat{\vec y}}
\def\hvz{\hat{\vec z}}
\def\hvX{\hat{\vec X}}

%------------------------------ Sets ---------------------------------
\def\SA{\mathcal A}
\def\SB{\mathcal B}
\def\SE{\mathcal E}
\def\SF{\mathcal F}
\def\SG{\mathcal G}
\def\SI{\mathcal I}
\def\SL{\mathcal L}
\def\SM{\mathcal M}
\def\SN{\mathcal N}
\def\SO{\mathcal O}
\def\SP{\mathcal P}
\def\SS{\mathcal S}
\def\SU{\mathcal U}
\def\SV{\mathcal V}
\def\SW{\mathcal W}
\def\SX{\mathcal X}
\def\SY{\mathcal Y}
\def\SZ{\mathcal Z}

%------------------------------ Misc ---------------------------------
\def\ii{^{(i)}}
\def\ij{^{(j)}}
\def\G{G(\veta, \gamma)}

%--------------------------- Densities -------------------------------
\def\fc{\tilde f(y|x)}
\def\fvc{\tilde f(\vy|\vx)}
\def\fcxy{f(x|y)}
\def\fcvxy{f(\vx|\vy)}
\def\fcyx{f(y|x)}
\def\fcvyx{f(\vy|\vx)}
\def\fija{f_{ij}(a)}
\def\fijak{f_{ij}^k(a)}
\def\fh{f(y, x; \veta)}
\def\fvh{f(\vy, \vx; \veta)}
\def\fkex{f_k^e(x_k)}
\def\fkevx{f_k^e(\vx_k)}
\def\fkpx{f_{k}^p(x_{k})}
\def\fkkpx{f_{k+1}^p(x_{k+1})}
\def\fkpvx{f_{k}^p(\vx_{k})}
\def\fkkpvx{f_{k+1}^p(\vx_{k+1})}
\def\ftxa{f(x_{k+1}|x_k, a_k)}
\def\ftvxa{f(\vx_{k+1}|\vx_k, \va_k)}
\def\fkv{f_k^v(v_k)}
\def\fkvv{f_k^v(\vv_k)}
\def\fkw{f_k^w(w_k)}
\def\fkvw{f_k^w(\vw_k)}
\def\fnex{f_n^e(x_n)}
\def\fnevx{f_n^e(\vx_n)}
\def\fnpx{f_{n}^p(x_{n})}
\def\fnnpx{f_{n+1}^p(x_{n+1})}
\def\fnpvx{f_{n}^p(\vx_{n})}
\def\fnnpx{f_{n+1}^p(\vx_{n+1})}
\def\fx{f(x)}
\def\fxy{f(x,y)}
\def\fvx{f(\vx)}
\def\fvxy{f(\vx,\vy)}
\def\fy{f(y)}
\def\fvy{f(\vy)}

%--------------------------- Distributions ---------------------------
\def\Fc{\tilde{F}(y|x)}
\def\Fh{F(y, x; \veta)}

%------------------------------ Theorems ------------------------------------
\theoremheaderfont{\normalfont\bfseries\sffamily}
\theoremstyle{plain}
{\theorembodyfont{\itshape} \newtheorem{theorem}     {Satz}[chapter]}
{\theorembodyfont{\itshape} \newtheorem{lemma}       {Lemma}[chapter]}
{\theorembodyfont{\itshape}\newtheorem{corollary}   {Korollar}[chapter]}
%{\theorembodyfont{\sffamily}\newtheorem{example}     {Beispiel}[chapter]}
{\theorembodyfont{\normalfont} \newtheorem{remark}      {\sffamily Remark}[chapter]}
{\theorembodyfont{\itshape}\newtheorem{definition}  {Definition}[chapter]}
{\theorembodyfont{\itshape}\newtheorem{axiom}  {Axiom}[chapter]}
%\theoremheaderfont{\scshape\normalsize}
%\theoremstyle{plain}
%{\theorembodyfont{\normalfont\small} \newtheorem{Proof}	{Proof}}
\newenvironment{proof}
{\noindent{\scshape Beweis.}}
{\hspace*{\fill}$\square$\bigskip}

\makeatletter 
\def\ifEmpty#1{\def\@temp{#1}\ifx\@temp\@empty} 
\makeatother

\newcounter{example}[chapter]
\renewcommand{\theexample}{\thechapter.\arabic{example}}
\newenvironment{example}[1][]
{\refstepcounter{example}\par\bigskip\noindent%
   {\normalfont\bfseries\sffamily Beispiel~\theexample~\ifEmpty{#1}\else(#1)\fi} \sffamily}
{\hspace*{\fill}$\blacksquare$\smallskip\smallskip}

\newenvironment{principle}[1][]
{\par\bigskip\noindent%
   {\normalfont\bfseries\sffamily \ifEmpty{#1}\else#1\fi} \itshape}
{\smallskip\smallskip}
\algrenewcommand{\algorithmiccomment}[1]{\hfill// #1}

\newcommand{\emphc}[1]{\comment{\normalfont\marginfont #1}\emph{#1}}

% Hervorheben & Index
\newcommand{\emphi}[2][]{%
	\emph{#2}%
	\ifthenelse{\isempty{#1}}{\index{#2}}{\index{#1}}%
}

% Hervorheben, Index & Marginalkommentar
\newcommand{\emphci}[2][]{%
	\comment{\normalfont\marginfont #2}%
	\emph{#2}%
	\ifthenelse{\isempty{#1}}{\index{#2}}{\index{#1}}%
}

%\input{macros/TableCommands}
\usetikzlibrary{topaths}
\usetikzlibrary{decorations.pathmorphing}
\usetikzlibrary{fit}

\newcommand{\RA}{\tikz[baseline=0mm]  \draw[ultra thick, -stealth] (0,1.3mm)--(3.5mm,1.3mm); }

\newcommand{\coordcross}[6]{%
	\draw[->, >=latex] (#1,0) -- (#2,0) node[right,below] {#5};
  \draw[->, >=latex] (0,#3) -- (0,#4) node[left] {#6};
}

%\newcommand{\xtick}[3]{%
%}

\tikzstyle{blockdef}=[rectangle, rounded corners, draw=black, thick, text centered]
\tikzstyle{rblock}=[blockdef, inner sep=6pt, text width=2.7cm]
\tikzstyle{cblock}=[blockdef, circle, inner sep=2pt, text width=2mm]
\tikzstyle{ablock}=[ultra thick, -latex, black]
\tikzstyle{aablock}=[ultra thick, latex-latex, black]
\tikzstyle{fblock}=[draw=black, fill]
\tikzstyle{fitblock}=[blockdef, fill=none, dashed, inner sep = 1mm]


%% Color Boxes
\tikzstyle{hlNone}=[anchor=base, rounded corners=3pt]
\tikzstyle{hlFrameNone}=[anchor=base, rounded corners=3pt, thick]
\tikzstyle{hlFrameBlue}=[hlFrameNone, fill=BlockBlue, draw=HilightBlue]
\tikzstyle{hlFrameRed}=[hlFrameNone, fill=BlockRed, draw=HilightRed]
\tikzstyle{hlBlue}=[hlNone, fill=BlockBlue]
\tikzstyle{hlGreen}=[hlNone, fill=BlockGreen]
\tikzstyle{hlLightGreen}=[hlNone, fill=BlockGreen!40]
\tikzstyle{hlRed}=[hlNone, fill=BlockRed]

\everymath{\displaystyle}
\def\hlFrameBlue<#1>#2{\alt<#1>{
  \tikz[baseline]{\node[hlFrameBlue] {#2};}
}{
  \tikz[baseline]{\node[hlFrameNone] {#2};}
}}

\def\hlFrameRed<#1>#2{\alt<#1>{
  \tikz[baseline]{\node[hlFrameRed] {#2};}
}{
  \tikz[baseline]{\node[hlFrameNone] {#2};}
}}

\def\hlBoxBlue<#1>#2{\alt<#1>{
  \tikz[baseline]{\node[hlBlue] {#2};}
}{
  \tikz[baseline]{\node[hlFrameNone] {#2};}
}}


\def\hlBoxGreen<#1>#2{\alt<#1>{
  \tikz[baseline]{\node[hlGreen] {#2};}
}{
  \tikz[baseline]{\node[hlFrameNone] {#2};}
}}

\def\hlBoxLightGreen<#1>#2{\alt<#1>{
  \tikz[baseline]{\node[hlLightGreen] {#2};}
}{
  \tikz[baseline]{\node[hlFrameNone] {#2};}
}}

\def\hlBoxRed<#1>#2{\alt<#1>{
  \tikz[baseline]{\node[hlRed] {#2};}
}{
  \tikz[baseline]{\node[hlFrameNone] {#2};}
}}
%%% Silbentrennung
\hyphenation{Strategie-iteration}
\hyphenation{Werte-iteration}
\hyphenation{Sicherheits-�quivalenz}

%%% Wasserzeichen "Entwurf"
\draft{%
\AddToShipoutPicture{%
  \setlength{\unitlength}{1cm}
  %\ifthenelse{\equal{\value{page}}{1}}{}{%
  \ifthenelse{\isodd{\value{page}}}
  	{\put(19,14.98){% Mitte rechts von A4 (21cm max 29.7cm)
      \makebox(0,0){\rotatebox{90}{\textcolor[gray]{0.8}{\fontsize{4cm}{4cm}\selectfont{\rm{Entwurf}}}}}
    }}
    {\put(2,14.98){ % Mitte links von A4 (21cm max 29.7cm)
        	\makebox(0,0){\rotatebox{-90}{\textcolor[gray]{0.8}{\fontsize{4cm}{4cm}\selectfont{\rm{Entwurf}}}}}
    }}
  }
%}
}


%%% Formatierung
%\clubpenalty=1000000
%\widowpenalty=1000000
%\brokenpenalty=1000000
%\finalhyphendemerits=1000000

%% Dokument Beginn %%%%%%%%%%%%%%%%%%%%%%%%%%%%%%%%%%%%%%%%%%%%%%%%%%%%%%%%

% - Deckblatt,
% - Inhaltsverzeichnis,
% - Hauptteil gegliedert z.B. in
%   Einleitung, Grundlagen, Experimente, Ergebnisse, Zusammenfassung
% - Literaturverzeichnis,
% - Abbildungsverzeichnis (ggf.),
% - Tabellenverzeichnis (ggf.),
% - Abkürzungsverzeichnis (ggf.),
% - Formelverzeichnis (ggf.),
% - Anhang, (nicht mehr Bestandteil der Arbeit! Wird daher nicht bewertet)
% - Erklärung der Urheberschaft,

\title{My Bachelor Thesis}
\author{Marcel Pascal Stolin}
\matrikel{32168	}
\thesis{Bachelor Thesis}
\geburtsdatum{03.04.1993}
\geburtsort{Kamen}
\abgabedatum{\today}
\supervisor{M.Sc. Christoph Hennebold}

\zusammenfassung{Hier kommt eine deutschsprachige Zusammenfassung hin.}
\abstract{Abstract in English.}

\begin{document}
% Deckblatt
\begin{titlepage}
\begin{minipage}[h]{1.0\linewidth+1.6cm}
  \begin{center}
		\rule[2ex]{\textwidth}{1pt}
		\begin{tabular}{ccc}
		 \includegraphics[width=.4\textwidth]{logo_iff} & \hphantom{abcdefghijklmnop} & \includegraphics[width=.3\textwidth]{logo_ipa}
		\end{tabular}
		\\[3mm]
    \rule[10ex]{\textwidth}{1pt}
		\huge\textbf{\textsf{\thethesis}}\\[2cm]
  	\Large \textbf{\textsc{\thetitle}} \\[1cm]
    \normalsize
    vorgelegt von\\[1ex]
		{\large \textbf{\theauthor}}\\[1ex]
		\thematrikel \\[1ex]
		geboren am \thegeburtsdatum~in~\thegeburtsort \\
    \vspace{2.5cm}
    Verfasst am\\[1ex]
		\textbf{\textsf{Fraunhofer-Institut f�r Produktionstechnik und Automatisierung IPA}}\\[1ex]
		und\\[1ex]
		\textbf{\textsf{Institut f�r Industrielle Fertigung und Fabrikbetrieb,\\Universit�t Stuttgart}}\\
    \vspace{2cm}
  \end{center}
	\normalsize{
    \begin{tabular}{ll}
    	Pr�fer: & Univ.-Prof. Dr.-Ing. habil. Marco Huber \\
      Betreuung: & \thesupervisor \\
    \end{tabular}\\[2ex]
    }
	\rule[-2ex]{\textwidth}{1pt}
  \begin{flushright}
		\footnotesize Stuttgart, \theabgabedatum
	\end{flushright}
\end{minipage}
\end{titlepage}


\frontmatter

\cleardoublepage
\chapter*{Eidesstattliche Erkl�rung}

Ich versichere hiermit, dass ich, \theauthor, die vorliegende \thethesis\ selbstst�ndig angefertigt, keine anderen als die angegebenen Hilfsmittel benutzt und sowohl w�rtliche als auch sinngem�� entlehnte Stellen als solche kenntlich gemacht habe. Die Arbeit hat in gleicher oder �hnlicher Form noch keiner anderen Pr�fungsbeh�rde vorgelegen.

\bigskip\bigskip\bigskip
\noindent Ort, Datum: \underline{\hphantom{01.01.2020, L�denscheidt-West}}

\bigskip\bigskip\bigskip
\noindent Unterschrift: $\underset{\text{\small\theauthor}}{\underline{\hphantom{\text{01.01.2020, L�denscheidt-West}}}}$

\cleardoublepage
\chapter*{Zusammenfassung}
\thezusammenfassung
%
\bigskip\bigskip\bigskip\bigskip\bigskip
\begin{flushright}
	{\LARGE\sffamily \textbf{Abstract}}
\end{flushright}
\bigskip\bigskip\bigskip
\theabstract

\cleardoublepage
% Inhaltsverzeichnis in den PDF-Links eintragen
% \pdfbookmark[1]{Inhaltsverzeichnis}{toc}
\tableofcontents

\chapter*{Notation}

\section*{Konventionen}
%
\begin{tabular}{rl}
$x$ & Skalar
\\
$\vx$ & Spaltenvektor
\\
$\x, \rvx$ & Zufallsvariable/-vektor
\\
$\hx, \hvx$ & Mittelwert/-vektor
\\
$x^*, \vx^*$ & Optimaler Wert/Vektor
\\
$x_{0:k}, \vx_{0:k}$ & Folge von Werten $(x_0, x_1, \ldots, x_k)$ / Vektoren $(\vx_0, \vx_1, \ldots, \vx_k)$
\\
$\A$ & Matrizen
\\
$\SA$ & Mengen
\\
$\preceq, \prec$ & schwache/strenge Pr�ferenzrelation
\\
$\NewR$ & Reelle Zahlen
\\
$\NewN$ & Nat�rliche Zahlen
\\
$\blacksquare$ & Ende eines Beispiels
\\
$\square$	& Ende eines Beweises
\end{tabular}

\section*{Operatoren}
%
\begin{tabular}{rl}
$\mat A\T$ & Matrixtransposition
\\
$\mat A^{-1}$ & Matrixinversion
\\
$|\mat A|$ & Determinante einer Matrix
\\
$|\SA|$ & Kardinalit�t der Menge $\SA$
\\
$\pot(\SA)$ & Potenzmenge von $\SA$
\\
$\E\{\cdot\}$ & Erwartungswertoperator
\\
$\SO(g)$ & O-Kalk�l entsprechend der Landau-Notation bei welcher \\&beispielsweise $f(x) \in \SO(g(x))$ besagt, dass die Funktion \\&$f(x)$ die Komplexit�t $\SO(g(x))$ besitzt
\end{tabular}

\section*{Spezielle Funktionen}
%
\begin{longtable}{rl}
$\Pr(\SE)$ & Wahrscheinlichkeitsma�, welches die Wahrscheinlichkeit
\\ & angibt, dass Ereignis $\SE$ eintritt
\\
$p(\vx)$ & (Wahrscheinlichkeits-)Dichtefunktion f�r kontinuierliche $\vx$ \\
& und Z�hldichte f�r diskrete $\vx$
\\
$p(\vx|\vy)$ & Bedingte Dichtefunktion
\\
$P(\vx)$ & (Wahrscheinlichkeits-)Verteilungsfunktion
\\
$\erf(x)$ & Gau�'sche Fehlerfunktion
\\
$\exp(x)$ & Exponentialfunktion $e^x$
\\
$\Gauss(\vx; \hvx, \C_x)$ & Gau�dichte, \dh Dichtefunktion eines normalverteilten \\&Zufallsvektors $\rvx$ mit Mittelwertvektor $\hvx$ und\\& Kovarianzmatrix $\C_x$
\end{longtable}

% Hauptteil
\mainmatter

%\iffalse
\chapter{Introduction}
\label{sec:demo}
%


%Complex computation operations rely on a IT infrastructure that is able to perform operations on scale. Computing systems that automatically adapt to demands and conditions of the workload serve the needs of
%highly scaled computing systems QUELLE. The microservice architecture is a technique to divide a big application into independent services. Each single service only serves one business task of the application QUELLE. In recent years, container applications have become popular to deploy microservices and visualize the needed resources for performing complex operations QUELLE. A container is a single unit that serves the 
%whole environment for a microservice.

% -- Struktur von Intro
% Grid computing
% Autonomic computing
% Microservices ??
% Container
% GPU acceleration


% Start
Storing huge amounts of data has become inexpensive in recent years, but processing it, requires 
parallel computations on clusters with multiple machines \cite{Chambers2018Spark}. Complex 
computation operations rely on a IT infrastructure with the ability to perform operations on scale. 
%Scalability is the ability of a system to grow in accordance of the amount of work \cite{TakadaDistr}.

% Sachen aus autonomic computing
In order to achieve high scalability, computing systems need to adapt dynamically to demands and conditions of the workload.


% Technologies
%In recent years, technologies have contributed to improve the scalability and performance of modern 
%computing systems like cloud computing, virtualization and GPU acceleration.

% Autonomic computing
%Autonomic computing is an architecture for computing systems to manage themselves in accordance to high level objectives, %configured by administrators. These computing systems dynamically adapt to demands and conditions of the workload. The f%our main attributes of autonomic computing are self-configuring, self-healing, self-optimizing and self-protecting \%cite{Kephart2003VisionComputing}. The autonomic computing architecture consists 
%of an autonomic manager, which is responsible to make adjustments in it`s own environment \cite{Sinreich2006AnAB}. %Technologies like Docker help to create containers ... bla bla bla

% Virtualization
%Modern virtualization technologies like containers help to overcome scalability limitations with the creation of virtual instances %of IT resources \cite{Erl2013CloudComputing}.


% GPU acceleration
% In addition to elasticity
%Another aspect to improve complex computation operations is GPU acceleration. GPU accelerations has been an on-going %research area in the recent years.

\section{Problem statement}

%% Struktur
% Möglichkeit IT resourcen automatisch nach bedarf zu skalieren
% Viele ETL operation sind daten intensiv von natur, GPU kann helfen

%To prevent a breakdown of the IT environment while complex operations are being computed, the IT environment should have the ability to scale IT resources dynamically
%when performance peaks and fluctuations occur.

% Dynamisch auf performance spitzen reagieren
ETL \footnote{Extract transform load} operations are compute-intensive. During the execution of analytic applications, performance thresholds can be reached and
the computing system can become out-of-order. 

In addition, many algorithms profit from data-parallelism. % Das ha mehr mit SPark zu tun

%If performance thresholds are reached and the IT environment can`t adapt 
%dynamically to the usage demand, the environment can look temporarily out-of-order. 
%During the execution of operations, performance threshold can be reached and leads to 

% Mit GPU ETLs beschleunigen
%In addition, algorithms do profit from data-parallelism. 

\section{Research questions}

\section{Thesis structure}
\chapter{Background}
\label{sec:background}
%
In this chapter bla bla
\todo{Describe Chapter}

% ====
% Autonomic computing
% ====
\section{Autonomic computing}
Autonomic computing is an architecture for computing systems to enable the ability to manage themselves in accordance to high level objectives configured by administrators \cite{Kephart2003VisionComputing}. 
These computing systems dynamically adapt to demands and conditions of the workload \cite{Kephart2003VisionComputing}.
An intelligent control-loop is responsible to collect all important details of the computing system and make decision according to the collected details. To automate the tasks, the intelligent control-loop is organized into four categories:

\begin{itemize}
\item \textbf{Self-configuring}
- Components in the environment have to adapt dynamically to system changes using policies. For example, deploying or removing new components.

\item \textbf{Self-healing}
- If system errors have being detected, the control-loop has to perform policy-based actions without disrupting the 
environment.

\item \textbf{Self-optimizing}
- The control-loop has to monitor the resources and should adapt to changes dynamically.

\item \textbf{Self-protecting}
- Detection and protection against threats.

\end{itemize}

 An autonomic computing environment consists of an autonomic manager, managed-resources and a knowledge-base.
 % As shown in figure ...
 
\subsection{Managed resources}
Managed resources are software or hardware components in the computing environment. For example, a managed resource 
can be a database, service, application, server or different entity. Each managed resource implements an interface to enable 
the autonomic manager to communicate with the managed resource. 

These interface are called touchpoints.

\subsection{Autonomic manager}
The autonomic manager implements an intelligent control-loop to collect system metrics from the managed resources and acts according to the collected details. It can only make adjustments within it`s own scope and uses policies to make decisions of what actions have to be
executed to accommodate the objectives.
To be self-managing, the autonomic manager has to implement the following four automated functions.

\begin{itemize}
\item \textbf{Monitor}
- The monitor function is responsible to collect the needed metrics from all managed resources and applies aggregation and filter
operations to the collected data. After that the function reports the metrics.

\item \textbf{Analyze}
- To determine if changes have to be made to the computing system, the collected data has to be analyzed.

\item \textbf{Plan}
- If changes have to be made, an appropriate change plan has to be generated. A change plan consists of actions that are needed to achieve the configured goals and objectives. The change plan needs to be forwarded to the execute function.

\item \textbf{Execute}
- The execute function applies all necessary changes to the computing system.

\end{itemize}

Multiple autonomic manager can exist in an autonomic computing environment to perform only certain parts. For example, 
there can be one autonomic manager which is responsible to monitor and analyze the system and another autonomic manager 
to plan and execute. To create a complete and closed control-loop, multiple autonomic manager can be composed together.


% ====
% Docker
% ====
\section{Docker}


% ====
% Apache Spark
% ====
\section{Apache Spark}
Apache Spark is an open-source computing framework for parallel data processing on a large computer cluster. Spark manages the available resources and distributes computation tasks across a cluster to perform big-data processing operations at large scale \cite{Chambers2018Spark}. Before Spark was developed, Hadoop MapReduce \cite{Dean2010MapReduce} was the framework of choice for parallel operations on a computer cluster \cite{Zaharia2010Spark}. Spark accomplished to outperform Hadoop by 10x for iterative Machine Learning \cite{Zaharia2010Spark}. It is implemented in Scala\footnote{Scala programming language. https://www.scala-lang.org/}, a JVM-based language and provides a programming interface for Scala, Java\footnote{Java programming language. https://www.oracle.com/java/}, Python\footnote{Python programming language. https://www.python.org/} and R\footnote{R programming language. https://www.r-project.org/}. In addition, Spark includes an interactive SQL shell and libraries to implement Machine Learning and streaming applications \cite{Chambers2018Spark}.
It was developed in 2009 as the Spark research project at UC Berkeley and became an Apache Software Foundation project in 2013 \cite{Chambers2018Spark}. 


\subsection{Spark programming model}
\label{sec:spark_programming_model}
% Explain RDDs
Spark provides resilient distributed datasets (RDDs) as the main abstractation for parallel operations \cite{Zaharia2010Spark}. Core types of Spark`s higher-level structured API are built on top of RDDs \cite{Chambers2018Spark} and will automatically be optimized by Spark`s Catalyst optimizer to run operations quick and efficient \cite{Hien2018Spark}.
\todo{Master-Slave Architektur + Bild}


\paragraph{Resilient distributed datasets:}
% What are RDDs
Resilient distributed datasets are fault-tolerant, parallel data structures to enable data sharing across cluster applications \cite{Zaharia2012RDDs}. They allow to express different cluster programming models like MapReduce, SQL and batched stream processing \cite{Zaharia2012RDDs}. RDDs have been implemented in Spark and serve as the underlying data structure for higher level APIs (Spark structured API) \cite{Zaharia2012RDDs}.
% How can RDDs be used in applications    !! HIER MAL BEISPIEL WIE IN RDD ERZEUGT WIRD
RDD`s are a immutable, partitioned collection of records and can only be initiated through transformations (e.g. map, filter) on data or other RDD`s.
% What are the advantages of RDDs    !!! HIER NOCHMAL BEISPIEL WIE DAS AUSSEHEN KANN
An advantage of RDDs is, that they can be recovered through lineage. Lost partitions of an RDD can be recomputed from other RDDs in parallel on different nodes \cite{Zaharia2012RDDs}. 
% Better use structured API
RDDs are lower level APIs and should only be used in applications if custom data partitioning is needed \cite{Chambers2018Spark}. It is recommended to use Sparks structured API objects instead. Optimizations for RDDs have to be implemented manually while Spark automatically optimize the execution for structured API operations \cite{Chambers2018Spark}.


\paragraph{Spark structured API:}
% What is the structured API
Spark provides high level structured APIs for manipulating all kinds of data. The three distributed core types are Datasets, DataFrames and SQL Tables and Views \cite{Chambers2018Spark}.
% About DataFrames, Datasets and SQL
Datasets and DataFrames are immutable, lazy evaluated collections that provide execution plans for operations \cite{Chambers2018Spark}. SQL Tables and Views work the same way as DataFrames, except that SQL is used as the interface instead of using the DataFrame programming interface \cite{Chambers2018Spark}.
% Differences between Datasets ad Dataframes
Datasets use JVM types and are therefore only available for JVM based languages. DataFrames are Datasets of type Row, which is the Spark internal optimized format for computations. This has advantages over JVM types which comes with garbage collection and object instantiation \cite{Chambers2018Spark}.


\paragraph{Spark Catalyst:}
% What is the Catalyst optimizer
Spark also provides a query optimizer engine called Spark Catalyst. \Fig{fig:spark_catalyst_process} illustrates how the Spark Catalyst optimizer automatically optimizes Spark applications to run quickly and efficient.
% Short How
Before executing the user`s code, the Catalyst optimizer translates the data-processing logic into a logical plan and optimizes the plan using heuristics \cite{Hien2018Spark}. After that, the Catalyst optimizer converts the logical plan into a physical plan to create code that can be executed \cite{Hien2018Spark}.


% What are logical plans
Logical plans get created from a DataFrame or a SQL query. A logical plan represents the data-processing logic as a tree of operators and expressions where the Catalyst optimizer can apply sets of rule-based and cost-based optimizations \cite{Hien2018Spark}.
For example, the Catalyst can position a filter transformation in front of a join operation \cite{Hien2018Spark}.

% What are physical plans
From the logical plan, the Catalyst optimizer creates one ore more physical plans which consist of RDD operations \cite{Chambers2018Spark}. The cheapest physical will be generated into Java bytecode for execution across the cluster \cite{Hien2018Spark}.

% Spark Catalyst figure
\begin{figure}[h]%
\centering
\includegraphics[scale=0.5]{images/03_background/spark_catalyst}%
\caption{Optimization process of the Spark Catalyst - Source: Authors own model, based on \cite{Hien2018Spark}.}%
\label{fig:spark_catalyst_process}%
\end{figure}
\todo{Bild nochmal machen mit abstand	+ QUELLE anpassen}


\subsection{Spark application architecture}

\begin{figure}[h]%
\centering
\includegraphics[scale=0.5]{images/03_background/cluster_overview}%
\caption{Overview of a Spark cluster architecture - Source: Authors own model, based on \cite{Apache2020Spark}.}%
\label{fig:spark_cluster_overview}%
\end{figure}

\todo{Lieber Doch Master/Slave ???. Was ist mit dem CLuster-Manager}
% Explain image
\Fig{fig:spark_cluster_overview} illustrates the main architecture of a Spark cluster. The architecture follows the master-worker model where the Spark driver is the master and the Spark executors are the worker \cite{Hien2018Spark}.
\todo{Nicht ganz richtig, Master und Worker sind machines und driver und executor sind prozesse}

\paragraph{Spark driver:}
The Spark driver is a JVM process on a physical machine and responsible to maintain the execution of a Spark application \cite{Chambers2018Spark}. It coordinates the application tasks onto each available executor \cite{Hien2018Spark}. To get launch executors and get physical resources, the Spark driver interacts with the cluster manager \cite{Chambers2018Spark, Hien2018Spark}.


\paragraph{Spark Executor:}
A Spark executor performs the tasks given by the Spark driver \cite{Chambers2018Spark}. It runs as a JVM process and runs until the Spark application finishes \cite{Hien2018Spark}. After the executor finishes, it reports back to the Spark driver \cite{Chambers2018Spark}. Each task will be performed on a separate CPU core to enable parallel processing  \cite{Hien2018Spark}.


\paragraph{Cluster manager:}
The cluster manager is an external service that orchestrates the work between the machines in the cluster \cite{Hien2018Spark, Apache2020Spark}. The cluster manager knows about the resources of each worker and decides on which machine the Spark driver process and the executor processes run \cite{Hien2018Spark, Chambers2018Spark}.
% Different cluster types
Spark supports different services that can run as the cluster manager: Standalone, Apache Mesos\footnote{Apache Mesos. https://mesos.apache.org/}, Hadoop YARN\cite{Murthy2013Yarn} and Kubernetes\footnote{Kubernetes. https://kubernetes.io/} \cite{Apache2020Spark}.
% What is standalone
\todo{Explain standalone}
% Explain different cluster deploy modes
The cluster manager provides three different deploy modes for acquiring resources in the cluster.
\begin{itemize}
\item Cluster mode
\item Client mode
\item Local mode
\end{itemize}

% Cluster mode
To run an application in cluster mode, the user has to submit a precompiled JAR, python script or R script to the cluster manager \cite{Chambers2018Spark}. After that, the cluster manager starts the driver process and executor processes exclusively for the Spark application on machines inside the cluster \cite{Chambers2018Spark, Hien2018Spark}.
\begin{figure}[h]%
\centering
\includegraphics[scale=0.5]{images/03_background/cluster_mode}%
\caption{Spark`s cluster mode - Source: Authors own model, based on \cite{Chambers2018Spark}.}%
\label{fig:spark_cluster_mode}%
\end{figure}

% Client mode
The difference between the client mode and the cluster mode is that, the driver process runs on the client machine outside of the Spark cluster \cite{Chambers2018Spark}.
\begin{figure}[h]%
\centering
\includegraphics[scale=0.5]{images/03_background/client_mode}%
\caption{Spark`s client mode - Source: Authors own model, based on \cite{Chambers2018Spark}.}%
\label{fig:spark_client_mode}%
\end{figure}

% Local mode
The local mode starts a Spark application on a single computer \cite{Chambers2018Spark}. It is not recommended to use the local mode in production, instead it should be used for testing Spark applications or learning the Spark framework \cite{Chambers2018Spark}.


\subsection{Spark application implementation}
% Short intro
The concept of a Spark application consists of calling transformations and actions. A transformation creates a DataFrame or a Dataset, the logical data structures of a Spark application. The computation of a Spark application gets processed when an action gets called in the application. The transformations of a Spark application build up a directed acyclic graph (DAG) of instructions. By calling an action, the DAG will break down into stages and tasks to create a single job for execution \cite{Chambers2018Spark}.
% Show example application
\begin{lstlisting}[frame=single, label=lst:spark_python_example, caption=Example of a Python3 Spark application, captionpos=b]
# Initialize a SparkSession
sparkSession = SparkSession\
    .builder\
    .getOrCreate()

# Create a dataframe with a transformation
dataframe = sparkSession.range(1, 1000)
# Apply another transformation
dataframe = dataframe.filter(dataframe.id % 2 == 0)
# Call an action
count = dataframe.count()
\end{lstlisting}
% Explain a spark application on top of an example
\Lst{lst:spark_python_example} demonstrates an example implementation of a Spark application. At first a SparkSession gets initialized. Each Spark application must include a SparkSession to initialize the application driver and executors  \cite{Chambers2018Spark}. In Addition, the SparkSession provides an API for data-processing logic and configuration of the Spark application \cite{Hien2018Spark}. After that, a DataFrame gets created with the range transformation to include each number from 1 to 1000 in the DataFrame. Next, a filter transformation is applied on the DataFrame to sort out any odd number. At the end, the number of rows gets saved in a variable with the count action.
\todo{Create listing macro}



% How to submit
\todo{Tabelle von Luu18 vll als Anhang(Table 2-1 Subdirectories Inside the spark-2.1.1-bin-hadoop2.7 Directory)}
The Spark framework provides a spark-submit executable to launch a Spark application inside a cluster.
% Show example submit
\begin{lstlisting}[frame=single, label=lst:spark_submit_example, caption=Execution of a Spark Python application using the spark-submit executable, captionpos=b]
$SPARK_HOME/bin/spark-submit \
    --master spark://spark-master:7077 \
    application.py
\end{lstlisting}
% Explain
\Lst{lst:spark_submit_example} provides an example how the spark-submit executable can be used to launch a Spark Python application.


\subsection{Spark standalone cluster deployment}
% Why only standalone
\todo{Why only standalone}
% What is a standalone cluster
The standalone mode is a basic cluster-manager build specifically for Spark. It is developed to only run Spark but supports workloads at large scale \cite{Chambers2018Spark}.

% How can we create a cluster
Spark provides build-in scripts to start a master node and worker nodes in standalone mode. ABC demonstrates how a master node and worker node gets launched using the build-in scripts.
% NE ALTER !, schon erwähnen start-master und start-worker


% ====
% NVIDIA RAPIDS
% ====
\section{RAPIDS accelerator for Apache Spark}
% Explain
The RAPIDS accelerator for Apache Spark is a plugin suite to enable GPU acceleration for computing operations on Apache Spark 3.x \cite{SparkRapids2020Docs}. 
% What it does
To accelerate computing operations, it uses the RAPIDS\footnote{Open GPU Data Science - https://rapids.ai/} libraries \cite{SparkRapids2020Docs} and extends the Spark programming model (see \SubSec{sec:spark_programming_model}) \cite{Mcdonald2020SparkRapids}.

\subsection{Extension of the Spark programming model}
% How does it extends DataFrame and SQL
The plugin suite extends the Spark programming model with a new DataFrame based on Apache Arrow\cite{ApacheArrow2020Docs} data structures and the Catalyst optimizer to generate GPU-aware query plans \cite{Mcdonald2020SparkRapids}.


% What is Apache Arrow
Apache arrow is a data platform to build high performance applications that work with large dataset`s and to improve analytic algorithms. A component of Apache Arrow is the Arrow Columnar Format, an in-memory data structure specification for efficient analytic operations on GPUs and CPUs \cite{ApacheArrow2020Docs}.


% How do DataFrames and SQL use RAPIDS
Spark`s DataFrame and SQL use the RAPIDS APIs to run tranformations and actions on a GPU.
% How does the Catalyst create GPU aware plans
The Spark Catalyst optimizer identifies operator in a query plan that are supported by the RAPIDS APIs. To execute the query plan, these operators can be scheduled on a GPU within the Spark cluster \cite{Mcdonald2020SparkRapids}. If operators are not supported by the RAPIDS APIs, a physical plan for CPUs will be generated by the Catalyst optimizer to execute RDD operations \cite{Mcdonald2020SparkRapids}. \Fig{fig:rapids_query_plan} illustrates, how a query plan gets optimized with the RAPIDS accelerator for Spark enabled.

% RAPIDS query plan image
\begin{figure}[h]%
\centering
\includegraphics[scale=0.25]{images/03_background/rapids_query_plan}%
\caption{Catalyst optimization with RAPIDS accelerator for Apache Spark - Source: Authors own model, based on \cite{Mcdonald2020SparkRapids}.}%
\label{fig:rapids_query_plan}%
\end{figure}


\subsection{GPU aware scheduling}


% ====
% Prometheus
% ====
\section{Prometheus}
% What is Prometheus
Prometheus is an open-source monitoring and alerting system \cite{Prom2020Docs}.
% Time-series format + PromQL
To collect and store data, Prometheus supports a multi-dimensional key-value pair based data model which can be analyzed in real-time using the PromQL query language \cite{Pandey2020Monitoring}.
% Pull-based approach
It follows the pull-based approach to scrape metrics from hosts and services \cite{Bastos2019Prom}.

% Pull-based image
\begin{figure}[h]%
\centering
\includegraphics[scale=0.25]{images/03_background/prom_pull-based}%
\caption{Pull-based approach to scrape metrics - Source: \cite{Bastos2019Prom}.}%
\label{fig:prom_pull-based}%
\end{figure}

% Describe figure
As \Fig{fig:prom_pull-based} demonstrates, a pull-based monitoring system scrapes metrics from services which makes them available for the monitoring system. In this case, the monitoring system needs a list of all hosts and services to monitor \cite{Bastos2019Prom}.


% What is monitoring and why is it important
A monitoring system keeps track of the health status of components in a computing cluster. Because of the complexity of modern computing clusters, the effort is to high to observe components manually \cite{Bastos2019Prom}.


% What are metrics
A metric is an observed property of software or hardware, e.g. the usage of a CPU core. To measure metrics, data-point`s will be recorded at a fixed time interval. A data-point is a pair of a value and a timestamp. The combination of multiple data-point`s is called a time-series \cite{Turnbull2016Monitoring}.


\subsection{Prometheus architecture}

\todo{Bilder für components}

% Architecture image
\begin{figure}[h]%
\centering
\includegraphics[scale=0.25]{images/03_background/prom_architecture}%
\caption{Prometheus high-level architecture - Source: Authors own model, based on \cite{Prom2020Docs, Bastos2019Prom}.}%
\label{fig:prom_architecture}%
\end{figure}

% High level architecture
\Fig{fig:prom_architecture} illustrates the high-level architecture of Prometheus.
% Description of components
The Prometheus ecosystem provides multiple components. Components can be optional, depending on the monitoring needs of the environment \cite{Bastos2019Prom}. The main components of a Prometheus system are Prometheus server, Alertmanager, service discovery, exporters, Pushgateway and visualization tools \cite{Prom2020Docs}.


\paragraph{Prometheus server:} 
The Prometheus server is the main component of a Prometheus system. It is responsible to collect metrics as time-series data from targets and stores the collected data in the built-in TSDB \cite{Bastos2019Prom}. Prometheus uses the concept of scraping to collect metrics from a target. A target host has to expose an endpoint to make metrics available in the Prometheus data format \cite{Pandey2020Monitoring}. Additionally, the Prometheus server triggers alerts to the Alertmanager if a configured condition becomes true \cite{Prom2020Docs}.


\paragraph{Alertmanager:}
If an alerting rule becomes true, the Prometheus server generates an alert and pushes it to the Alertmanager. The Alertmanager generates notifications from the received alerts. A notification can take multiple forms like emails or chat messages. Webhooks can be implemented to trigger custom notifications \cite{Bastos2019Prom}.


\paragraph{Service discovery:}
As mentioned before, Prometheus follows a pull-based approach to fetch metrics from a target. To know about all targets, Prometheus needs a list of the corresponding hosts. The service discovery manages the complexity of maintaining a list of hosts manually in an changing infrastructure \cite{Bastos2019Prom}.


\paragraph{Exporters:}
If an application does not support an endpoint for Prometheus, an exporter can be used to fetch metrics and make them available to the Prometheus server. An exporter is a monitoring agent running on a target host that fetches metric from the host and exports them to the Prometheus server \cite{Pandey2020Monitoring}.


\paragraph{Pushgateway:}
If a target is not designed to be scraped, metrics can be pushed against the Pushgateway\cite{Prom2020Docs}. The Pushgateway converts the data into the Prometheus data format and passes them to the Prometheus server \cite{Pandey2020Monitoring}.


\paragraph{Visualization:}
Prometheus supports various tools for virtualization of the scraped data. Grafana\footnote{Grafana: The open observability platform - https://grafana.com/} is one of the widely used tools for this occasion.


%\subsection{Creating alerts}
% Beschreiben wie man das mach in der yml datei


\subsection{Monitoring Docker container}
% How can Prometheus be used with Docker (cAdvisor)
% Auch auf andere Technologien eingehen -> Measuring DockerPerformance What a Mess!!!
\todo{Wahrscheinlich erst nach der implementation}


\section{Gitlab CI/CD}


\section{K-MEANS}


\section{Naive Bayes Classifier}


\section{Scaling heat}


\section{KHP}
\chapter{Related Work}
\label{sec:related}
%

This chapter provides background information about ...
\todo{Describe Chapter}


% ===========================================
% ===========================================
% Technologies
% ===========================================
% ===========================================
\section{State-of-the-art technologies}


% ===========================================
% ===========================================
% Scalable container
% ===========================================
% ===========================================
\section{Scalable container architectures}
% State of the art monitoring
Dynamic scaling of containerized application is an active area of research. In recent years containerized applications have been used efficiently and many research works focused on creating a infrastructure with state-of-the-art monitoring tools for measuring container performance and the ability to dynamically scale. %Other work published about auto-scaling policies


% Delivering elastic ...
Barna et al. \cite{Barna2017ElasticContainerApps} proposed a development approach for an autonomic management system for containerized applications enabling DevOps\footnote{DEVOPS BLA BLA}.
% The goal
The goal of this work was to monitor the performance of the system under low or high workload by using an autonomic management architecture to create a model to determine necessary scaling actions. They proposed the Scaling Heat Algorithm to determine if the container environment has to perform a scale-out or scale-in action. In addition DEVOPS
% Results
It was observed, that the resulting model makes accurate scaling decisions with an error rate less than 10\%.


% Automated Deployment of a Spark Cluster with Machine Learning Algorithm Integration
Fernández et al. \cite{Fernandez2020AutomatedDeploySpark} introduced LadonSpark, a solution to automate the deployment process of an Apache Spark cluster.
% The goal
The objective of LadonSpark is a method for deploying and configuring an Apache Spark Cluster as a task.


% Application deployment using containers with auto-scaling for microservices in cloud environment
Srirama et al. \citep{Srirama2020AppDeplyCont} designed a heuristic-based auto-scaling strategy for container-based microservices in a cloud environment.
% The goal
The purpose of the auto-scaling strategy was to balance the overall resource utilization across microservices in the environment.
% Results
The proposed auto-scaling strategy performed better results than state-of-the-art algorithms in processing time, processing cost and resource utilization. The processing cost of microservices could be reduced by 12-20\% and the CPU and memory utilization of cloud-servers have been maximized by 9-15\% and 10-18\%.


% ===========================================
% ===========================================
% GPU Stuff
% ===========================================
% ===========================================
\section{Heterogenous GPU aware Spark systems}
% Why is GPU needed
Apache Spark is a computing framework that distributes tasks on multiple CPU cores. Data and compute intensive applications profit from GPU acceleration. Therefore, various research projects took effort to bring GPU acceleration to Apache Spark.


% HeteroSpark
Li et al. \cite{Li2015HeteroSpark} created a GPU-accelerated heterogeneous architecture called HeteroSpark to integrate with Apache Spark.
% The goal
HeteroSpark contributes to enable data parallelism for data and compute intensive applications running on Apache Spark. To be able to enable or disable GPU acceleration, HeteroSpark provides a plug-n-play design.
% Results
Overall, HeteroSpark is able to achieve a 18x speed-up for various Machine Learning applications running on Apache Spark.


% HetSpark
Klodjan et al. \cite{Klodjan2018HetSpark} introduced HetSpark a heterogeneous modification of Apache Spark.
% The goal
HetSpark extends Apache Spark with two executors, a GPU accelerated executor and a commodity class. The GPU accelerated executor is based on VineTalk\cite{Mavridis2017VineTalk} for GPU acceleration.
% Results
The authors observed, that for compute intensive tasks GPU accelerated executers are preferable while for linear tasks CPU-only accelerators should be used.


% Spark-GPU
Yuan et al. \cite{Yuan2016SparkGPU} proposed SparkGPU to enable parallel processing with GPUs in Apache Spark and contributes to achieve high performance and high throughput in Apache Spark applications.
SparkGPU extends Apache Sparks to determine the suitability of parallel-processing for a task to enable task scheduling between CPU and GPU. 
SparkGPU accomplished to improve the performance of machine learning algorithms up to 16.13x and SQL query execution performance up to 4.83x.



\chapter{Conceptual Design}
\label{sec:design}
%
\todo{Describe Chapter}


% ===========================================
% ===========================================
\section{Cluster Architecture}

\begin{figure}[h]%
\centering
\includegraphics[scale=0.8]{images/04_conceptual_design/cluster_architecture/overall_architecture}%
\caption{Scaling process UML activity model - Source: Authors own model.}%
\label{fig:ca-overall_architecture}%
\end{figure}

% Cluster consists of
The cluster conists of multiple components runnin in a own docker container. Prometheus, cAdvisor, Spark Master, Spark Worker, Auto-Scaler.
% Communication between components
cAdvisor collects metrics from all available docker container running in the same network. Prometheus scrapes metrics from cAdvisor and stores them as time-series-data in its ts-db. The Auto-Scaler read metric from Prometheus via its REST API. If a scaling action is necessary the auto-scaler scales the number of worker.
% 


% ===========================================
% ===========================================
\section{DevOps Process}


% ===========================================
% ===========================================
\section{Autonomic Manager}
% Short description of old section
As described in section ABC, the Autonomic Manager follows the MAPE Architecture. The Autonim manager implements a Loops which is responsible for collecting and analyzing the metrics.


\subsection{Workflow}
% Process figure
\begin{figure}[h]%
\centering
\includegraphics[scale=0.6]{images/04_conceptual_design/01_scaling_process_model}%
\caption{Scaling process UML activity model - Source: Authors own model.}%
\label{fig:am-workflow-process_model}%
\end{figure}

% Whats is the workflow/meaning

% How is the workflow structured
As illustrated in \Fig{fig:am-workflow-process_model} the workflow consists of four different steps:

\begin{enumerate}
\item Collect metrics
\item Analyze metrics
\item Create scaling plan
\item Scale environment
\end{enumerate}

In the first step, all metrics (described in \SubSec{subsec:design-metrics}) are collected using cAdvisor and will be stored in Prometheus. The autonomic manager (described in \Sec{sec:design-manager}) gathers all metrics from Prometheus REST API and analyzes the data to discover if scaling actions are necessary. According to the data, the manager creates a plan to estimate how many nodes needed to be added or removed. In the last step, the scaling plan will be executed by the manager.

% Wie werden Metriken gesammelt
% Wie werden Metriken analysiert
% Wei wird der Manager benachrichtigt
% Wie wird skaliert
% Die Loop implementiert die loop von autonomous computing MAPE


\subsection{Design}

% Component design
\begin{figure}[h]%
\centering
\includegraphics[scale=0.6]{images/04_conceptual_design/autonomic_manager/design}%
\caption{Scaling process UML activity model - Source: Authors own model.}%
\label{fig:am-design-component}%
\end{figure}

The autonomic manager will be compose by several components to build a manager. To collect the metrics, cAdvisor will be used (see Section A), for storing the metrics Prometheus will be used (see Section B) and for auto-scaling, a Auto-scaler will be developed.


\section{Auto-Scaler}
% How is the autonimic manager configurated
% What components does the auto-scaler consist of? Configuration, Analyze, Plan

\subsection{Configuration}
The Auto-Scaler needs specific configuration properties to be able to collect the correct metric from Prometheus and deploy new Apache Spark worker container in the environment. The following are properties that have to be defined to ensure that the Auto-Scaler is able to collect meaningful metrics and scale Apache Spark worker as expected.

\subsubsection{General properties}

\begin{itemize}
\item \textbf{Interval seconds:} The number of seconds when the loop has to repeat needs to be defined.
\end{itemize}

\subsubsection{Apache Spark worker properties}

\begin{itemize}
\item \textbf{Worker image:} To guarantee that each Spark worker is homogeneous, all worker container should be created with the same image.

\item \textbf{Worker network:} To establish communication between all Spark worker and the Spark master, all new Spark worker container should be in the same network.

\item \textbf{Worker thresholds:} The minimum and maximum number of concurrent Spark worker should be defined. To avoid the cold start effect, the minimum amount of worker should be 1. \todo{CHeck nochmal den Cold Start effekt}

\item \textbf{Apache Spark master URI:} To distribute the workload across all Spark Worker, all Spark Worker need to communicate with the Spark master. The Spark master URI starts with \textit{spark://}.
\end{itemize}

\subsubsection{Prometheus properties}

\begin{itemize}
\item \textbf{Prometheus URL:} The Auto-Scaler will collect the configured metrics from Prometheus REST API.

\item \textbf{Prometheus query:} A PromQL query needs to be defined to collect specific metrics for CPU utilization calculation. 
\end{itemize}


\subsection{Analyze}
% Hier scaling heat algorithmus bla bla

\subsection{Plan}
% Hier KHPA Algorithmus

\subsection{Scaling}
% There are threshold defined over and under utilization
% There are min and max thresholds of worker defined

If the CPU Utilization goes above 80\%, it means the cluster is overloaded. If the Utilization is less than 20\% the cluster is underutilized.


% ===========================================
% ===========================================
\section{Metrics}
% Welche metrics gibt es
% Welche metrics werden gesammelt (nur von den Worker)
% Wie werden die metrics berechnet

\subsection{CPU Utilization}
% Evtl ist hier der faösche platz -> Besser eine Metrics section wo erklärt wird, wie das berechnet werden soll.
\todo{Hier auf Tabelle verweisen (Anhang) Metriken die von Prometheus + cAdvisor bereitgestellt werden}

% How is the CPU Utilization being calculated
To adapt to business needs, the CPU percentage of each Spark Worker will be calculated. Prometheus provides several metrics to calculate the CPU percentage. The CPU percentage of all Worker can be calculated as follows:

% Ist nicht ganz korrekt, muss ja irgendwie noch durch die Anzahl an Worker geteilt werden
\begin{equation}
CPUUtilization = \sum ContainerCPUUserSecondsTotal \times 100 / NumberOfWorker
\label{eq:formel}
\end{equation}


\subsection{GPU Utilization}


\chapter{Implementation}
\label{sec:implementation}
%
\todo{Describe Chapter}

\section{Cluster architecture}


\chapter{Evaluation}
\label{sec:evaluation}
%
\todo{Describe Chapter}

\section{Cluster architecture}


\chapter{Outlook}
\label{sec:outlook}
%
\todo{Describe Chapter}

\section{Cluster architecture}


\chapter{Conclusion}
\label{sec:conclusion}
%
\todo{Describe Chapter}

\section{Cluster architecture}


%\fi


% Anhang
%\appendix
% 'Anhang' ins Inhaltsverzeichnis
%\phantomsection
%\addcontentsline{toc}{part}{Anhang}
%\input{content/appendix_}

% Anhang (Bibliographie darf im deutschen nicht in den Anhang!)
\bibliography{literatur}
%\nocite{*}
\clearpage

% Abbildungs- und Tabellenverzeichnis
\listoffigures
%\listoftables


\IfDefined{printindex}{\printindex}
\IfDefined{printnomenclature}{\printnomenclature}



%% Dokument ENDE %%%%%%%%%%%%%%%%%%%%%%%%%%%%%%%%%%%%%%%%%%%%%%%%%%%%%%%%%%
\end{document}

